% =====================================================================
%  Phase C --- Visual alignment on the painted mark
%  SCOPE: the downward camera, the detector, the bounded correction.
% =====================================================================
\section{Closing the Last Centimetres}
\label{sec:phaseC_align}

\subsection{Why odometry is not enough here}

Sec.~\ref{sec:phaseC_catalogue} established the constraint: two stations
sharing an inner aisle are \SI{0.10}{m} apart. Sec.~\ref{sec:slam}
established the measurement: uncalibrated odometry on this platform
diverges past \SI{10}{m} over a mission, and even calibrated it holds
only \measured{\SI{0.20}{m}}.

\SI{0.20}{m} is twice the station spacing. The failure this produces is
not an imprecise reading --- it is a \emph{correctly measured value
filed under the wrong label}, which is worse, because nothing downstream
can detect it. A campaign so corrupted looks complete.

\subsection{The sensor point is the camera, not the base}

\phaseC{} adds a downward camera on a short boom over the front bumper,
and makes \emph{that} point the one a station names.

A robot that has a body cannot see the floor under its own middle --- the
chassis is in the way --- so a mark under the centre can be driven to but
never verified. Putting the reference point under a downward camera makes
``the mark is in the middle of the picture'' and ``the sensor is on the
mark'' the same statement, with no lever arm whose sign can be got wrong.

\SI{0.50}{m} is the smallest offset that also keeps the camera mast out
of the downward view cone: the mast reaches $x = 0.27$~m at \SI{0.45}{m}
height, and a camera at \SI{0.50}{m} sees the floor from $x = 0.27$~m
outwards, so no ray crosses it. Two consequences follow and are stated
once: the base centre stands \SI{0.50}{m} back along the aisle from the
station, and \textbf{every pose in every report is the sensor point},
not \texttt{base\_link}.

\subsection{Detecting a cross, and refusing a blob}

The detector (\filename{agri/vision.py}) is a red-dominance threshold
followed by a connected-component pass, with three refusals that matter
more than the detection itself:

\begin{itemize}[leftmargin=*,itemsep=1pt,topsep=2pt]
  \item \textbf{A solid blob is not a cross.} A cross fills
        \measured{$\approx$44\%} of its own bounding box; a filled square
        fills 100\%. Rejecting on fill ratio stops the robot from
        centring on a red object that is not a station mark. The
        accepted band is deliberately wide, because the exact ratio
        moves with the arm and thickness constants, and a detector that
        needed retuning whenever the paint changed would be retuned
        wrong.
  \item \textbf{A clipped cross is refused, not averaged.} A mark half
        out of frame has a centroid that is not its centre. Averaging it
        yields a confident correction in the wrong direction.
  \item \textbf{The nearer of two marks wins.} With a neighbour
        \SI{0.10}{m} away the frame can contain both. The suite asserts
        the nearer one is selected at three different offsets, including
        one where the robot starts closer to the neighbour.
\end{itemize}

The threshold itself needs no tuning: the aisle paint is dark green
$(38, 64, 54)$ and the gutters are near white, so a factor-1.6 dominance
test separates the marker from both by a wide margin. This is a
consequence of the world being generated (Sec.~\ref{sec:phaseC_catalogue})
--- the detector and the paint come from the same file.

\subsection{The correction is bounded, and the bound is not a preference}

The visual correction is capped at \measured{\SI{0.04}{m}} per station.

The reasoning is the same \SI{0.10}{m} that drove the arm length: a
correction larger than \SI{0.05}{m} could leave the robot nearer the
\emph{neighbouring} station than the one it is about to file under. The
Cloud would then reject the report --- correctly --- for a reason with
nothing to do with the plant. \SI{0.04}{m} keeps a clear margin under
that bound and is still four times the residual the odometric approach
leaves behind.

The robot makes at most four correction attempts, stopping when the
residual falls below \SI{0.008}{m}.

\subsection{What happens when the camera cannot see}

The mark can be missing, occluded, or outside the frame. \phaseC{} does
not treat that as a mission failure and does not silently pretend it
did not happen. The robot parks on odometry alone, says so in the log,
and \textbf{increments a counter that travels with the mission report}:

\begin{lstlisting}[style=bash]
mission 788081147cfd finished in 601 s, floor camera
used at 48 station(s), unavailable at 0
\end{lstlisting}

That line is the honest form of the accuracy claim. Rather than
asserting a parking error the simulator cannot independently confirm,
the system reports \emph{how many of the visits were closed visually}
and lets the reader judge. In the reference campaign of
Sec.~\ref{sec:results} the camera resolved its mark at
\measured{48 of 48} stations, and at \measured{52 of 52} across the full
session including the four single-station re-measurements.

Every report additionally carries its own \param{parking\_error\_m}, and
the Cloud checks it against the catalogue independently of what the
robot believed --- so a visit filed under the wrong label is detectable
after the fact, from the stored record alone.
